\begin{topic}[id=ota_strategies,
             difficulty={Mittel},
             type={Architekturvergleich},
             prerequisites={Grundlagen eingebetteter Systeme; Basiswissen Software-Updates und Security hilfreich},
             practice={optional},
             owner={Dozententeam},
             updated={2026-04-09},
             tags={ OTA, Uptane, SUIT, Signaturen, Rollback, Root of Trust, Policy }]{OTA-Update-Strategien}

\TopicDescription{Sichere Updates sind zentral für vernetzte Geräte. Dieses Thema gibt eine Landkarte: Paketierung/Signierung, sichere Boot-Kette, Geräte-Identität, Versionierung/Kompatibilität und Rollout-Strategien (Canary, Staged, Rollback). Du lernst, was Uptane und IETF-SUIT liefern, wo kleine Geräte an Grenzen stoßen (Flash-Budget, A/B-Slot vs. Single-Image) und wie Telemetrie/Monitoring in den Prozess passt.}

\TopicLeitfragen{\item Welche Kriterien entscheiden zwischen SUIT und Uptane?
\item Wie wird Rollback sicher und nutzbar umgesetzt?
\item Welche Telemetrie braucht ein Update-System?
}
\TopicScope{\item Keine vollständige End-to-End-Implementierung.
\item Kein Hardware-Bootloader-Coding.
}
\TopicPractice{Sequenzdiagramm für Download, Verifikation und Rollback.}

\TopicKeywords{OTA, Uptane, SUIT, Signaturen, Rollback, Root of Trust, Policy}

\nocite{uptaneSpec,ietfSUIT}
\end{topic}
